\section{Thesis Workspace and Supporting Documents}
\label{sec:thesis-workspace}

\subsection{A separate writing deliverable}

The root \path{thesis/} directory is an English LaTeX skeleton, not a finished
thesis or another detector implementation. The small \path{thesis.tex} entry
file specifies document order. Personal details and the approved title are
centralised in \path{metadata.tex}; layout and PDF metadata have their own
files. This report does not reproduce the student's personal metadata.

Each front-matter item has a separate source: frontispiece, English abstract,
Italian summary, optional acknowledgements, contents, figure list, table list,
and acronyms. Seven chapter files contain short prompts and headings for the
introduction, background, implementation, experimental method, results,
discussion, and conclusions. The bibliography and reproducibility appendix are
also separate. This avoids a single source file containing the entire thesis.

\subsection{Template and build provenance}

The workspace adapts the community Overleaf DEI template by
Cecco di Ronchitti da Bruzene under CC BY 4.0. Its publicly displayed source is
preserved in \path{thesis/upstream/main.tex}. Exporting the complete original
Overleaf project ZIP requires an account; the local copy is explicitly
described as the public-source adaptation, not an unmodified complete ZIP.
The university seal comes unchanged from the official English frontispiece.
Template attribution and the changes are recorded in the workspace README.

\path{thesis/build.sh} uses the existing \path{compose.latex.yaml} service
with LuaLaTeX and latexmk. It works from the repository root or thesis folder,
and writes the PDF and auxiliary files below \path{out/thesis/}.
The original \path{scripts/latex.sh} remains available for the technical
reports using their established PDFLaTeX setup. The Compose service is shared;
the detector image does not acquire a TeX distribution.

The skeleton was compiled and its cover inspected when created. Metadata,
English labels, all front-matter items, seven chapter headings, and short
source files were checked. It uses 12-point text, 1.5 spacing, and the indicative
DEI margins. Times New Roman is selected if installed; the existing container
instead uses a disclosed TeX Gyre Termes draft fallback. PDF/A-2b settings and
metadata are configured, but the final document still needs actual conformance
validation. A metadata declaration alone is not certification.

\subsection{Length, writing guide, and scope}

The supervisor-approved project budget is a minimum of 60 main-text pages,
a target of 80--100, and a cap of 100. It counts the Introduction through the
Conclusions, including figures and tables, but excludes front matter,
bibliography, and appendices. This is the agreed writing budget, not a numerical
university page rule. The writing guide records the official sources checked
and the distinction between a file-size limit and a page count.

\path{docs/thesis_writing_guide.md} also records the chapter plan, citation and
visual conventions, evidence requirements, writing order, and final review.
This project report is a larger reference for those chapters; it is not
governed by the submitted thesis's chapter/page budget. Results should be
generated from retained artefacts, and text should be written in the student's
own academic voice under the applicable university guidance.

\subsection{How the older documents fit together}

The proposal describes the broad original architecture and research questions.
The status and adaptation reports describe earlier implementation boundaries.
The LM implementation plan, change record, row-wise meeting report, paper
comparison, and professor-style experiment record explain numerical decisions.
The overdispersion and window-length notes support the statistical discussion.
The temporal review and branch reports explain remaining data/evaluation work.
The edge guideline and deployment notes describe different stages of the
gateway deliverable.

Earlier status documents and implementation plans are archived design and
experiment records. Their original boundaries do not define an active software
backlog. This report and the completion record identify the delivered scope;
optional extensions are collected in Section~\ref{sec:future-work}.
